\chapter{Évaluation et gestion des hallucinations}

% Partie consacrée à la qualité du RAG et à la détection / prévention
% des hallucinations.

\section{Protocole d'évaluation}
% - Présentation du jeu de cas d'évaluation
%   (cf. \texttt{rag\_pi/evaluation\_cases.py}).
% - Types de cas : questions factuelles, questions hors-sujet, questions
%   pièges (refus attendu).
% - Métriques :
%   \begin{itemize}
%     \item présence des termes attendus dans la réponse,
%     \item reconnaissance des sources / articles cités,
%     \item détection des réponses de refus,
%     \item grounding (réponse appuyée sur les documents retournés).
%   \end{itemize}

\section{Mise en oeuvre}
% - Script \texttt{scripts/evaluate\_rag.py}.
% - Sortie : \texttt{evaluation/results/latest.json}
%   et \texttt{evaluation/results/latest.csv}.
% - Exécution reproductible.

\section{Résultats}
% - Tableau de résultats par type de cas.
% - Discussion : forces, points faibles, sensibilité aux paramètres
%   (\texttt{CHUNK\_SIZE}, \texttt{RETRIEVER\_K}, seuil de pertinence).

% Exemple de tableau à remplir :
%\begin{table}[ht]
%    \centering
%    \begin{tabular}{lcc}
%        \hline
%        Type de cas & Précision & Taux de refus correct \\
%        \hline
%        Factuels & \dots & \dots \\
%        Pièges / hors-sujet & \dots & \dots \\
%        \hline
%    \end{tabular}
%    \caption{Résultats d'évaluation par catégorie.}
%    \label{tab:eval}
%\end{table}

\section{Gestion des hallucinations}
% - Stratégies mises en place :
%   \begin{itemize}
%     \item seuil minimal de pertinence,
%     \item prompt forçant la citation des articles,
%     \item refus explicite quand aucun document pertinent n'est trouvé,
%     \item affichage des sources dans l'interface.
%   \end{itemize}
% - Cas concrets observés (avant / après).
